% !TeX root = ../../金融经济学_王江_讲义及习题.tex
% 08_组合选择一 课后习题（由章节文件拆分）

\begin{problem}
\problemno{8.1} 假设有 $N$ 项资产，收益率为 $\tilde{r}_{i}$，$i=1,\cdots,N$，它们是独立同分布的。对于一个具有不满足性和凹性效用函数的投资者，计算他的最优投资组合。
\end{problem}
\begin{solution}
设组合权重 $\theta=(\theta_{1},\ldots,\theta_{N})$，$\sum_{i}\theta_{i}=1$，期望效用目标先写为
\[
 \max_\theta E\!\left[u\!\left(
 \underbrace{w\sum_i\theta_i\tilde r_i}_{\tilde w\text{：期末财富}}
 \right)\right].
\]
其一阶条件为：对任意 $j,k$，
\[
    E\!\left[u'\!\left(w\sum_i\theta_i\tilde r_i\right)(\tilde r_j-\tilde r_k)\right]=0.
\]
由于 $\{\tilde{r}_{i}\}$ 独立同分布，目标函数对 $\theta$ 对称；严格凹效用下最优解唯一，故必有 $\theta_{i}^{*}=1/N$。即最优组合是所有风险资产的等权重组合。
\end{solution}

\begin{problem}
\problemno{8.2*} 假设你必须将 $w$ 投资于两只证券：一只无风险债券和一只风险证券。无风险债券的收益率确定为 $r_{F}$。风险证券的收益率为 $\tilde{r}$，它的均值为 $\bar{r}$ 而方差为 $\sigma^{2}$。假设你的目标就是最大化你的二次效用函数的期望值：
\[
    E\Bigl[\tilde{w}-\tfrac12 a\tilde{w}^{2}\Bigr].
\]
\begin{enumerate}[label=(\alph*)]
    \item 求解最优组合。
    \item 讨论最优组合如何依赖于股票和债券的期望收益率 $r_{F}$ 和 $\bar{r}$，股票收益率的波动率 $\sigma$，以及偏好系数 $a$。为这些依赖关系提供经济解释。
\end{enumerate}
\end{problem}
\begin{solution}
\begin{enumerate}[label=(\alph*)]
    \item 假设投资于股票的资产比例为 $x$，那么有
    \[
        \tilde{w}=w\bigl(1+r_{F}+x(\tilde{r}-r_{F})\bigr)
    \]
    参与者的优化问题是
    \[
        \max_x E\!\left[
 \underbrace{w\bigl(1+r_F+x(\tilde r-r_F)\bigr)}_{\tilde w}
 -\frac a2\bigl(w(1+r_F+x(\tilde r-r_F))\bigr)^2\right].
    \]
    先计算期末财富的一阶、二阶矩：
    \[
    \begin{aligned}
        E[\tilde w]&=w\bigl(1+r_F+x(\bar r-r_F)\bigr),\\
        E[\tilde w^2]&=w^2\bigl[(1+r_F+x(\bar r-r_F))^2+x^2\sigma^2\bigr].
    \end{aligned}
    \]
    代入目标函数后对 $x$ 求导，一阶条件为
    \[
        w(\bar r-r_F)-aw^2\bigl[(1+r_F+x(\bar r-r_F))(\bar r-r_F)+x\sigma^2\bigr]=0.
    \]
    整理得
    \[
        (\bar r-r_F)\bigl[1-aw(1+r_F)\bigr]
        =aw\bigl((\bar r-r_F)^2+\sigma^2\bigr)x.
    \]
    解得最优权重
    \[
        x^{*}=\frac{(\bar{r}-r_{F})\bigl(1-aw(1+r_{F})\bigr)}{aw\bigl(\sigma^{2}+(\bar{r}-r_{F})^{2}\bigr)}
    \]

    \item $x^{*}$ 随着 $\bar{r}$ 的增加而增加，随着 $r_{F}$、$\sigma^{2}$、$a$ 的增加而减小。股票的期望收益率增加使得股票的需求增加，而无风险收益率增加使得债券的需求增加；$\sigma^{2}$ 增大，股票的风险变大，对股票的需求变小；最后，$a$ 增大表明参与者更加厌恶风险，因而对股票的需求减小。
\end{enumerate}
\end{solution}

\begin{problem}
\problemno{8.3*} 考虑与练习 8.2 一样的问题。不过现在假设效用函数具有常数绝对风险厌恶系数：
\[
    E\bigl[-e^{-a\tilde{w}}\bigr].
\]
另外，假设股票收益率 $\tilde{r}$ 是正态分布的。
\begin{enumerate}[label=(\alph*)]
    \item 求解最优投资组合。
    \item 讨论最优组合如何依赖于股票和债券的期望收益率 $r_{F}$ 和 $\bar{r}$，股票收益率的波动率 $\sigma$，以及偏好系数 $a$。为这些依赖性提供经济解释。[如果需要的话，你可以添加其它条件。]
\end{enumerate}
\end{problem}
\begin{solution}
\begin{enumerate}[label=(\alph*)]
    \item 假设投资于股票的资产比例为 $x$，那么有
    \[
        \tilde{w}=w\bigl(1+r_{F}+x(\tilde{r}-r_{F})\bigr)
    \]
    参与者的优化问题是
    \[
        \max_x E\!\left[-\exp\!\left(-a
 \underbrace{w(1+r_F+x(\tilde r-r_F))}_{\tilde w}\right)\right].
    \]
    因 $\tilde{r}$ 正态，$\tilde{w}$ 亦正态。对正态随机变量有
    \[
        E\bigl[-e^{-a\tilde{w}}\bigr]=-\exp\Bigl(-aE[\tilde{w}]+\tfrac12 a^{2}\operatorname{Var}(\tilde{w})\Bigr)
    \]
    其中 $E[\tilde{w}]=w\bigl(1+r_{F}+x(\bar{r}-r_{F})\bigr)$，$\operatorname{Var}(\tilde{w})=w^{2}x^{2}\sigma^{2}$。最大化上式等价于最大化确定性等价，完整写为
    \[
        \max_x\left\{w\bigl(1+r_F+x(\bar r-r_F)\bigr)
        -\frac a2 w^2x^2\sigma^2\right\}.
    \]
    对 $x$ 的一阶条件为
    \[
        w(\bar{r}-r_{F})-a\,w^{2}x\sigma^{2}=0
    \]
    因而最优权重（CARA 下风险资产的金额 $wx^{*}$ 与财富无关）为
    \[
        x^{*}=\frac{\bar{r}-r_{F}}{aw\sigma^{2}}
    \]
    \item $x^{*}$ 随着 $\bar{r}$ 的增加而增加，随着 $r_{F}$、$\sigma^{2}$、$a$ 以及财富 $w$ 的增加而减小。超额收益上升提高股票需求；$\sigma^{2}$ 或绝对风险厌恶 $a$ 增大则降低股票需求。与二次效用不同：CARA 下风险资产金额 $wx^{*}=(\bar{r}-r_{F})/(a\sigma^{2})$ 不随 $w$ 变化，故权重 $x^{*}$ 随财富下降。
\end{enumerate}
\end{solution}

\begin{problem}
\problemno{8.4**} 赋予投资者初始金融财富 $w$ 以及 $\lambda$ 份不可交易收入 $\tilde{y}$，其中 $\lambda\ge 0$ 且 $E[\tilde{y}]=0$。他可以交易一只无风险证券和一只风险证券。无风险证券的收益率为 0 而风险证券的收益率 $\tilde{r}$ 的均值为 $\bar{r}$。他具有不满足性且是风险厌恶的。他的绝对风险厌恶系数随着他的财富增加而减少。令 $a$ 为风险证券上的投资额。他的期望效用函数为：
\[
    \max E\bigl[u(w+a\tilde{r}+\lambda\tilde{y})\bigr].
\]
\begin{enumerate}[label=(\alph*)]
    \item 当 $\lambda=0$ 时，风险资产的最优头寸 $a_{0}^{*}$ 如何随着 $w$ 变化而变化？
    \item 假设 $\tilde{y}$ 和 $\tilde{r}$ 是独立的。当 $\lambda>0$ 时，记风险资产的最优头寸为 $a_{\lambda}^{*}$。$a_{\lambda}^{*}$ 是大于还是小于 $a_{0}^{*}$？
    \item 假设 $\tilde{y}=c(\tilde{r}-\bar{r})+\tilde{\varepsilon}$ 以及 $E[\tilde{\varepsilon}\mid\tilde{r}]=0$。因此，如果 $c\neq 0$，$\tilde{y}$ 和 $\tilde{r}$ 是相关的。当 $\lambda$ 极小时，他持有的股票头寸如何随着 $w$ 变化而变化？解释所得到的结论。
\end{enumerate}
\end{problem}
\begin{solution}
\begin{enumerate}[label=(\alph*)]
    \item 当 $\lambda=0$ 时，优化的一阶条件是
    \[
        E\bigl[u'\!\left(\underbrace{w+a_0^*\tilde r}_{\tilde w}\right)\tilde{r}\bigr]=0
    \]
    两边对 $w$ 求导可得
    \[
        \frac{da_{0}^{*}}{dw}=\frac{E[u''(\tilde{w})\tilde{r}]}{-E[u''(\tilde{w})\tilde{r}^{2}]}
    \]
    这里绝对风险厌恶系数沿用正文的记号，写为
    \[
        \underbrace{-\frac{u''(t)}{u'(t)}}_{A(t)\text{：绝对风险厌恶系数}}.
    \]
    又
    \[
        \begin{aligned}
            E[u''(\tilde{w})\tilde{r}]
            &=-E[u'(\tilde{w})A(\tilde{w})\tilde{r}]\\
            &\ge -E[u'(\tilde{w})A(w)\tilde{r}]\\
            &=-A(w)E[u'(\tilde{w})\tilde{r}]\\
            &=0
        \end{aligned}
    \]
    故 $\dfrac{da_{0}^{*}}{dw}\ge 0$，即最优头寸随着 $w$ 的增加而增加；

    \item 参与者优化的一阶条件是
    \[
        E\bigl[u'\!\left(\underbrace{w+a_\lambda^*\tilde r+\lambda\tilde y}_{\tilde w}\right)\tilde{r}\bigr]=0
    \]
    两边对 $\lambda$ 求导可得
    \[
        \frac{da_{\lambda}^{*}}{d\lambda}=\frac{E[u''(\tilde{w})\tilde{r}\tilde{y}]}{-E[u''(\tilde{w})\tilde{r}^{2}]}
    \]
    将下面的条件期望记为 $K$：
    \[
        \underbrace{E[-A(\tilde w)\tilde y\mid\tilde r]}_{K(w+a_\lambda^*\tilde r)}.
    \]
    可以证明 $K(\cdot)>0$，$K'(\cdot)>0$，因此
    \[
        \begin{aligned}
 E[u''(\tilde w)\tilde r\tilde y]
 &=E[-A(\tilde w)u'(\tilde w)\tilde r\tilde y]\\
 &=E[K(w+a_\lambda^*\tilde r)u'(\tilde w)\tilde r]\\
 &\ge K(w)E[u'(\tilde w)\tilde r]=0.
 \end{aligned}
    \]
    最后一个等式是因为一阶条件，因此，$\dfrac{da_{\lambda}^{*}}{d\lambda}\ge 0$，故 $a_{\lambda}^{*}\ge a_{0}^{*}$；

    \item 由 (a) 有
    \[
        \frac{da_{0}^{*}}{dw}=\frac{E[u''(\tilde{w})\tilde{r}]}{-E[u''(\tilde{w})\tilde{r}^{2}]}=\frac{-E[u'(\tilde{w})A(\tilde{w})\tilde{r}]}{-E[u''(\tilde{w})\tilde{r}^{2}]}
    \]
    而
    \[
        \begin{aligned}
 E[u'(\tilde w)A(\tilde w)\tilde r]
 &=E[A(w+a\tilde r)u'(\tilde w)\tilde r]\\
 &\quad+\lambda E[A'(w+a\tilde r)u'(\tilde w)\tilde r\tilde y]
 +o(\lambda^2).
 \end{aligned}
    \]
    由 (a) 知道 $E\bigl[A(w+a\tilde{r})u'(\tilde{w})\tilde{r}\bigr]\le 0$，另外 $\lambda$ 又极小，因此
    \[
        E[u'(\tilde{w})A(\tilde{w})\tilde{r}]\le 0,\qquad \frac{da^{*}}{dw}\ge 0
    \]
    即最优头寸随着初始财富 $w$ 的增加而增加。这是因为 $\lambda$ 极小，因而不可交易收入不会改变 (a) 中的结论。
\end{enumerate}
\end{solution}

\begin{problem}
\problemno{8.5**} 假设你是对冲基金中交易政府债的交易员。你有 $w$ 的现金额度。
\begin{itemize}
    \item 在 0 期，当美联储拍卖新债券时，你可以提交报价而买入政府债。在 1 期，你可以在二级市场上以价格 $\tilde{v}$ 再把它们卖掉，其中 $\tilde{v}=\bar{v}+\tilde{\varepsilon}$，$\bar{v}>0$ 且 $E[\tilde{\varepsilon}]=0$。假设无风险收益率为零。
    \item 你购买债券的报价是 $p$。你购买到的数量 $x(p)$ 是 $p$ 的递增函数且 $x(0)=0$。选择 $p$ 以最大化你的期望效用函数：
    \[
        V(p)=E[u(\tilde{w})],\qquad \tilde{w}=w+x(p)(\bar{v}+\tilde{\varepsilon}-p).
    \]
    假设你是不满足的且严格风险厌恶的。为简单起见，假设 $V(p)$ 是严格凹的。
\end{itemize}
\begin{enumerate}[label=(\alph*)]
    \item 写出关于最优报价 $p^{*}$ 的一阶条件。
    \item 证明 $0<p^{*}<\bar{v}$。
    \item 对于两个绝对风险厌恶系数分别为 $A_{1}(w)$ 和 $A_{2}(w)$ 的交易员，如果 $\forall\,w$，$A_{1}(w)\ge A_{2}(w)$，那么 $p_{1}^{*}\le p_{2}^{*}$，其中 $p_{i}^{*}$ 是交易员 $i$ 的最优报价，$i=1,2$。
    \item 假设 $u(\cdot)$ 是绝对风险厌恶递减的。证明初始财富 $w$ 的增加使得 $p^{*}$ 也增加。
\end{enumerate}
\end{problem}
\begin{solution}
\begin{enumerate}[label=(\alph*)]
    \item 一阶条件是
    \[
        \begin{aligned}
 V'(p^*)&=E\!\left[
 u'\!\left(\underbrace{w+x(p^*)(\bar v+\tilde\varepsilon-p^*)}_{\tilde w}\right)\right.\\
 &\hspace{3em}\left.\cdot\bigl(x'(p^*)(\bar v+\tilde\varepsilon-p^*)-x(p^*)\bigr)\right]=0.
 \end{aligned}
    \]

    \item
    \[
        \begin{aligned}
            V'(0)&=E\bigl[u'(w)\bigl(x'(0)(\bar{v}+\tilde{\varepsilon})\bigr)\bigr]=u'(w)x'(0)\bar{v}>0\\
            V(\bar{v})&=E\bigl[u(w+x(\bar{v})\tilde{\varepsilon})\bigr]<u(w)=V(0)
        \end{aligned}
    \]
    又因为 $V(p)$ 是凹的，故 $0<p^{*}<\bar{v}$；

    \item 由定理 (7.4) 可知，存在 $f'>0$，$f''<0$ 使得 $u_{1}=f(u_{2})$。参与者的优化问题的一阶条件为：
    \[
        \begin{aligned}
 E\!\left[
 u_2'\!\left(\underbrace{w+x(p_2)(\bar v+\tilde\varepsilon-p_2)}_{\tilde w_2}\right)\right.&\\
 \left.\cdot\bigl(x'(p_2)(\bar v+\tilde\varepsilon-p_2)-x(p_2)\bigr)\right]&=0.
 \end{aligned}
    \]
    考虑
    \[
        V_{1}'(p_{2})=E\bigl[f'(u_{2}(\tilde{w}_{2}))u_{2}'(\tilde{w}_{2})\bigl(x'(p_{2})(\bar{v}+\tilde{\varepsilon}-p_{2})-x(p_{2})\bigr)\bigr]
    \]
    将报价的边际财富变化标记为
    \[
        \underbrace{x'(p_2)(\bar v+\tilde\varepsilon-p_2)-x(p_2)}_{\tilde r}.
    \]
    那么 $\tilde{w}_{2}$ 是 $\tilde{r}$ 的增函数，又 $f'(u_{2}(\tilde{w}_{2}))>0$ 且是 $\tilde{w}_{2}$ 的减函数，因此，$V_{1}'(p_{2})\le 0$，另外 $V(p)$ 是凹函数，因此 $p_{1}^{*}\le p_{2}^{*}$。

    \item 假设 $w_{1}>w_{0}$，当 $w=w_{0}$ 时优化问题的一阶条件为
    \[
        \begin{aligned}
 E\!\left[u'\!\left(\underbrace{w_0+x(p_0)(\bar v+\tilde\varepsilon-p_0)}_{\tilde w_0}\right)\right.&\\
 \left.\cdot\bigl(x'(p_0)(\bar v+\tilde\varepsilon-p_0)-x(p_0)\bigr)\right]&=0.
 \end{aligned}
    \]
    以下仍将绝对风险厌恶系数写为
    \[
        \underbrace{-\frac{u''(t)}{u'(t)}}_{A(t)}.
    \]
    等式两边对 $w_{0}$ 求导可得
    \[
        \begin{aligned}
            &E\bigl[u''(\tilde{w}_{0})\bigl(x'(p_{0})(\bar{v}+\tilde{\varepsilon}-p_{0})-x(p_{0})\bigr)\bigr]\\
            &=E\bigl[-A(\tilde{w}_{0})u'(\tilde{w}_{0})\bigl(x'(p_{0})(\bar{v}+\tilde{\varepsilon}-p_{0})-x(p_{0})\bigr)\bigr]\\
            &\ge 0
        \end{aligned}
    \]
    因此，
    \[
        E\bigl[u'\bigl(w_{1}+x(p_{0})(\bar{v}+\tilde{\varepsilon}-p_{0})\bigr)\bigl(x'(p_{0})(\bar{v}+\tilde{\varepsilon}-p_{0})-x(p_{0})\bigr)\bigr]\ge 0
    \]
    而 $V(p)$ 又是凹的，因此 $p_{1}\ge p_{0}$，即初始财富增加报价也会增加。
\end{enumerate}
\end{solution}

